\section{Intertemporal provision}
\label{sec:provision}

A medium that cannot be held must still let a member provide for later. Two instruments
answer, both already on the ledger: the claim, and capacity. What is not offered is a
third, an undated, counterparty-free stock --- a \emph{balance}. A balance is a claim with
the debtor's name erased, and every mechanism here consumes the name: insurance is a
reserved path from named underwriters, recourse is substitution onto them, discharge is a
named debtor's delivery. Erasing the name removes not the risk a claim carries but the
parties who priced it (Property~\ref{prop:no-hoard}).

\subsection{The claim: provision to a date}

A claim is dated, and the date is a term of the trade: a maturity between
$\kConstMinMaturityEpochs$ and $\kConstMaxHorizonEpochs$ epochs --- decades, at the
genesis epoch length --- movable only forward, by extension on both signatures
(\S\ref{sec:medium}). Its amount does not decay: decay is a property of stakes, and the
reservation behind an insured claim floors the edges it runs over (\S\ref{sec:standing}).
Insured, within the horizon below, it survives its debtor: substitution makes the holder
whole in claims on the underwriters, one layer and no further (\S\ref{sec:recourse}).
Nothing bounds what a member may be owed --- the capacity gate sits on the debtor, because
debt spends the community's underwriting and being owed spends nothing --- so a creditor
position of any size is expressible. The ledger keeps the claims themselves, never a
signed net position.

Provision to a date is therefore available, and priced. For the life of the claim the
debtor's capacity stays reserved and the underwriters' supply stays drawn
(Rem.~\ref{rem:uw-duration}); the holder pays nothing while the two parties not holding
pay continuously, so a long maturity is a real concession, and its price has a place to
live, the face amount agreed at acceptance. The underwriters are asked once, collectively
and in advance: a claim is insured only within the \emph{insured horizon} of its
acceptance --- $\kConstInsuredHorizonEpochs$ epochs at genesis, a governed constant between
the maturity floor and the maximum horizon (\S\ref{sec:governance}) --- and a claim booked
or extended past it is the creditor's own, on the creditor's own signature. A claim meant
to run for decades is held as a series, settled and re-accepted at each horizon, each
acceptance priced against the community's backing at that time, which is where the
underwriters are asked again. At community scale the same fact is a budget: every
long-dated insured claim occupies flow current trade would otherwise be insured by, so
provision and trade compete for the same cut, visibly rather than through an aggregate
that silently inflates. A saver draws down the way anything is paid here: by buying from
their debtors, which nets first (\S\ref{sec:sale}); by their claims following trade, as
debtors' own sales route obligations onward (\S\ref{sec:medium}); by delivery, which
settles.

\subsection{Capacity: provision across seasons}

Capacity is the other instrument, and the partial one: a member who borrows and honours
accumulates standing that lets them draw on the community later with nothing in hand. It
decays unless renewed, so it provides across seasons rather than decades; it is not
transferable, so it cannot be given to a child or left in a will; and it is a claim on a
community's willingness, not on a stock of anything.

\subsection{The net saver's ledger}

A member who produces more than they consume accumulates claims --- others' obligations
to them --- and only claims: selling confers no standing (Rem.~\ref{rem:owed-not-sold}),
so their capacity grows only on the side where they buy and honour. That makes visible
what a currency balance conceals: one person's savings are other people's obligations,
worth what those people deliver. The risk is allocated as everywhere else --- an insured
claim is covered up to what each underwriter declared and no further, an uninsured one is
the creditor's own (Rem.~\ref{rem:uninsured-alone}).

\subsection{The horizon, and institutions}

No claim outlives the graph that prices it. At a horizon of decades the residual risk is
not the mechanism but the world: the debtor, the one layer of underwriters behind them
and the community itself must still be there, trading, at the date. That risk is
irreducible in any closed credit system, and a balance would carry it identically, with
the names, the insurance and the price removed. Provision across decades --- pensions,
illness --- therefore belongs to institutions: a pension society or mutual fund is a
member like any other, holding an account, declaring a supply and so standing behind
credit with a liability the ledger enforces, nameable as arbiter or beneficiary, measured
by the same cut. The protocol supplies that substrate; creating such institutions belongs
to the community, which should plan to serve this horizon with them and with assets held
outside the medium.
